\section{Step 54 Separate literals, end-of-block and lengths}\label{png-huff-s54}
Fill \texttt{huff\_classify} after the supplied literal/length decode. Values below 256 select the literal handler; 256 selects the existing block-completion path. For larger values, require a length code from 257 through 285 before indexing length tables. Reserved fixed symbols 286 and 287 must fail.

The fixed-literals input first returns zero and selects \texttt{emit\_literal}. That byte is the first row's filter selector, not a packed pixel. End-of-block finishes this block rather than this whole call; the final flag still decides whether another block follows.

\paragraph{Observe} Use \texttt{s54-1}; stop before \texttt{emit\_literal}. Inspect the state described above before continuing.

\Needspace{7\baselineskip}
\paragraph{Check} What does symbol 256 do?
\begin{itemize}
\item[A.] Writes byte zero to output.
\item[B.] Ends this block and follows its completion path.
\item[C.] Supplies a match distance.
\end{itemize}
